% 似地行星（综述）
% license CCBYSA3
% type Wiki

本文根据 CC-BY-SA 协议转载翻译自维基百科\href{https://en.wikipedia.org/wiki/Earth_analog}{相关文章}。

\begin{figure}[ht]
\centering
\includegraphics[width=8cm]{./figures/849251199efbb53d.png}
\caption{地球与金星的演化路径。金星一直是一个典型范例，用来说明一颗外观上类似地球的行星，其演化结果可能会与地球产生何等巨大的差异。} \label{fig_Eaan_1}
\end{figure}
类地行星（Earth analog），也称为“地球双胞胎”或“第二地球”，是指其环境条件与地球相似的行星或卫星。“类地行星”（Earth-like planet）这一术语也常被使用，但该术语有时泛指任何类地行星。

这一可能性对天体生物学家和天文学家而言尤为重要，其基本逻辑在于：一颗行星与地球越相似，它就越有可能具备维持复杂地外生命的能力。因此，这一主题长期以来一直是科学、哲学、科幻文学以及大众文化中的重要猜想和讨论对象。空间殖民以及“太空与人类存续”理念的倡导者，也长期以来一直在寻找可供定居的类地行星。在遥远的未来，人类甚至可能通过行星改造（terraforming）来人工制造一颗类地行星。

在对系外行星进行科学搜索与研究之前，这一可能性主要通过哲学和科幻作品来探讨。哲学家曾提出，由于宇宙尺度极其巨大，某处必然存在一颗与地球几乎完全相同的行星。“平庸原理”认为，像地球这样的行星在宇宙中应当十分常见；而“稀有地球假说”则认为此类行星极其罕见。迄今为止已发现的成千上万个系外行星系统，在整体结构上与太阳系有着显著差异，这在一定程度上支持了稀有地球假说。

2013 年 11 月 4 日，天文学家基于开普勒空间望远镜的数据报告称，在银河系内围绕类太阳恒星和红矮星的宜居带中，可能存在多达 400 亿颗地球大小的行星 $^{[1][2]}$。从统计意义上看，距离地球最近的此类行星可能位于约 12 光年之内 $^{[1][2]}$。2020 年 9 月，天文学家在已确认的 4000 多颗系外行星中，依据天体物理参数以及地球生命的自然演化历史，识别出了 24 颗“超级宜居行星”（即在宜居性方面优于地球的行星）候选者 $^{[3]}$。
\begin{figure}[ht]
\centering
\includegraphics[width=8cm]{./figures/22fbf76c2e490700.png}
\caption{一颗假想的宜居系外行星及其三颗天然卫星的艺术示意图。} \label{fig_Eaan_2}
\end{figure}
自20世纪90年代以来的科学发现，极大地影响了天体生物学领域的研究范围、行星宜居性模型，以及对地外智能生命（SETI）的探索。
\subsection{历史}
